\section{USAMO 2026}
        \begin{enumerate}
        	\item (\textit{USAMO 2026}) Let $n$ be an integer greater than $1$. For which real numbers $x$ is 

 \[\lfloor nx\rfloor-\sum_{k=1}^n\frac{\lfloor kx\rfloor}{k}\]
maximal, and what is the maximal value that this expression can take?
Note: $\lfloor z\rfloor$ denotes the greatest integer less than or equal to $z$.


 

	\item (\textit{USAMO 2026}) Annie is playing a game where she starts with a row of positive integers, written on a blackboard, each of which is a power of $2$. On each turn, she can erase two adjacent numbers and replace them with a power of $2$ that is greater than either of the erased numbers. This shortens the row of numbers, and she continues to take turns until only one number remains. Annie wins the game if the final remaining number is less than $4$ times the sum of the original numbers. Is it always possible for Annie to win, regardless of the starting row of numbers?


 

	\item (\textit{USAMO 2026}) Let $ABC$ be an acute scalene triangle with no angle equal to $60^\circ$. Let $\omega$ be the circumcircle of $ABC$. Let $\triangle_B$ be the equilateral triangle with three vertices on $\omega$, one of which is $B$. Let $\ell_B$ be the line through the two vertices of $\triangle_B$ other than $B$. Let $\triangle_C$ and $\ell_C$ be defined analogously. Let $Y$ be the intersection of $AC$ and $\ell_B$, and let $Z$ be the intersection of $AB$ and $\ell_C$.


 Let $N$ be the midpoint of minor arc $BC$ on $\omega$. Let $\mathcal{R}$ be the triangle formed by $\ell_B$, $\ell_C$, and the tangent to $\omega$ through $N$. Prove that the circumcircle of $AYZ$ and the incircle of $\mathcal{R}$ are tangent.


 

	\item (\textit{USAMO 2026}) A positive integer $n$ is called solitary if, for any nonnegative integers $a$ and $b$ such that $a+b=n$, either $a$ or $b$ contains the digit $1$ when written in base $10$. Determine, with proof, the number of solitary integers less than $10^{2026}$.


 

	\item (\textit{USAMO 2026}) Let $ABC$ be a triangle. Points $D$, $E$, and $F$ lie on sides $BC$, $CA$, and $AB$, respectively, such that

 \[\angle AFE=\angle BDF=\angle CED.\]
Let $O_A$, $O_B$, and $O_C$ be the circumcenters of triangles $AFE$, $BDF$, and $CED$, respectively. Let $M$, $N$, and $O$ be the circumcenters of triangles $ABC$, $DEF$, and $O_AO_BO_C$, respectively. Prove that $OM=ON$.


 

	\item (\textit{USAMO 2026}) Let $a$ and $b$ be positive integers such that $\varphi(ab+1)$ divides $a^2+b^2+1$. Prove that $a$ and $b$ are Fibonacci numbers.


 Note: Euler’s totient function $\varphi(n)$ is the number of positive integers less than
or equal to $n$ that are relatively prime to $n$.


 

\end{enumerate}